\documentclass[10pt,twocolumn,letterpaper]{article}

% --- Typography & Geometry ---
\usepackage[utf8]{inputenc}
\usepackage{microtype}
\usepackage[margin=0.75in,top=0.85in,bottom=0.85in]{geometry}

% --- Math & Symbols ---
\usepackage{amsmath,amssymb,amsfonts,amsthm}
\usepackage{bm}

% --- Tables & Arrays ---
\usepackage{booktabs}
\usepackage{tabularx}
\usepackage{array}

% --- Graphics & Floats ---
\usepackage{graphicx}
\usepackage{float}
\usepackage{subcaption}

% --- Code Listings ---
\usepackage{xcolor}
\usepackage{listings}

\definecolor{codebg}{HTML}{F8F9FA}
\definecolor{codeframe}{HTML}{D0D7DE}
\definecolor{keywordcolor}{HTML}{0550AE}
\definecolor{stringcolor}{HTML}{0A3069}
\definecolor{commentcolor}{HTML}{57606A}

\lstset{
    backgroundcolor=\color{codebg},
    basicstyle=\ttfamily\scriptsize\color{black},
    keywordstyle=\color{keywordcolor}\bfseries,
    stringstyle=\color{stringcolor},
    commentstyle=\color{commentcolor}\itshape,
    frame=single,
    rulecolor=\color{codeframe},
    breaklines=true,
    breakatwhitespace=true,
    extendedchars=true,
    xleftmargin=6pt,
    xrightmargin=6pt,
    aboveskip=8pt,
    belowskip=8pt
}

% --- Hyperlinks ---
\usepackage{hyperref}
\hypersetup{
    colorlinks=true,
    linkcolor=blue!70!black,
    citecolor=blue!70!black,
    urlcolor=blue!70!black
}

% --- Input Title & Authors ---
\section{00 Meta}

Draft content for nexus.


\begin{document}

\maketitle

% --- Abstract & Keywords Block ---
\begin{abstract}
We present \textbf{NEXUS} (\textbf{N}eural \textbf{E}fficient \textbf{X}-scoring \textbf{U}nified \textbf{S}parse 
attention), a modular attention framework that unifies six recent advances in 
sparse attention: (1) XAttention antidiagonal block scoring (training-free), 
(2) SeerAttention learnable gating, (3) DuSA dual-stage intra/inter-block 
sparsity, (4) AnchorAttention stripe band masking, (5) Spectraformer-style 
Random Fourier Feature global rectification, and (6) softpick rectified softmax.

The core mechanism partitions the sequence into blocks of size $B_s$, scores 
block-to-block interaction with coarse subsampled correlations, retains the 
top-$k$ key blocks per query block, and computes exact intra-block attention 
with optional stripe masking. A learnable Seer gating network provides a 
data-dependent override, and periodic RFF layers inject a global receptive 
field.

On SST-2, NEXUS achieves \textbf{{{D\_ACC}}\%} on DistilBERT (baseline {{BASE\_D}}\%, 
$\Delta$ = {{D\_DELTA}}), \textbf{{{R\_ACC}}\%} on RoBERTa (baseline {{BASE\_R}}\%, $\Delta$ = {{R\_DELTA}}), 
and \textbf{{{G\_ACC}}\%} on GPT-2 (baseline {{BASE\_G}}\%, $\Delta$ = {{G\_DELTA}}), using 
{{RECOVERY\_STEPS}} LoRA recovery steps.

\textbf{Important honest caveats:}

\begin{itemize}
\item \textbf{Encoder-only sparsity.} The block-sparse path is used only for
  non-causal models. Causal decoders (GPT-2) always fall back to dense
  attention, so NEXUS provides \textbf{no compute savings on GPT-2}.
\item \textbf{Stripe masking is gated at $N > 256$.} For SST-2 ($N = 128$), the stripe
  component is inactive.
\item \textbf{RFF rectification is disabled by default} when pretrained weights are
  provided. It activates only when \texttt{enable\_spectraformer=True}.
\item \textbf{Softpick activation} occasionally produces near-uniform distributions
  when all logits are negative; this is mitigated by our temperature
  parameter.

\end{itemize}
We present NEXUS as a \textbf{modular framework} for studying sparse attention 
combinations, not as a single breakthrough method.

\textbf{Keywords:} sparse attention, block routing, mixture of block attention, 
softpick, random Fourier features, multi-scale sparsity



\vspace{0.5em}
\noindent\textbf{Keywords:} sparse attention, block routing, mixture of block attention,
\end{abstract}

\vspace{1em}

% --- Paper Body ---
\section{Introduction}

\subsection{Motivation}

Sparse attention has seen rapid progress in 2024--2025. Notable recent 
advances include:

\begin{itemize}
\item \textbf{XAttention} (Xu et al., ICML 2025): training-free block scoring via
  antidiagonal correlations.
\item \textbf{SeerAttention} (Gao et al., NeurIPS 2025): learnable per-head block
  sparsity.
\item \textbf{DuSA} (Cao et al., NeurIPS 2025): dual-stage intra-block dense +
  inter-block sparse.
\item \textbf{AnchorAttention} (Zhu et al., 2025): stripe-based fine-grained sparsity.
\item \textbf{Spectraformer} (Nguyen et al., 2025): RFF-based global rectification.
\item \textbf{Softpick} (Zuhri et al., 2025): rectified softmax with natural ~47\%
  sparsity.

\end{itemize}
\textbf{These techniques are complementary but rarely combined.} NEXUS is our 
attempt to unify them in a single attention module.

\subsection{Problem Statement}

Given a pretrained transformer $\mathcal{M}$ with $H$ heads and hidden 
dimension $d$, we seek an attention mechanism $\mathcal{A}'$ such that:

\[
\text{FLOPs}(\mathcal{A}') \ll \mathcal{O}(N^2), \quad
\mathcal{A}' \approx \mathcal{A}
\tag{1.1}
\]

subject to:

\begin{enumerate}
\item \textbf{Modularity:} Each sparse mechanism can be enabled/disabled.
\item \textbf{Multi-architecture:} DistilBERT, RoBERTa, GPT-2.
\item \textbf{Pretrained compatibility:} Reuse pretrained Q/K/V/O weights.
\item \textbf{Causal correctness:} Preserve causality for decoders.
\item \textbf{Training stability:} No collapses.

\end{enumerate}
\subsection{Our Approach}

NEXUS combines six sparse-attention mechanisms in a unified pipeline:

\subsubsection{Block Partitioning}

The sequence is partitioned into $M = \lceil N / B_s \rceil$ blocks of size 
$B_s$ (default 16).

\subsubsection{XAttention Block Scoring}

Sub-sampled queries and keys compute coarse block-to-block scores:

\[
S_{b_i, b_j} = \frac{1}{|B_i| |B_j|} \sum_{i \in B_i, j \in B_j} \frac{q_i \cdot k_j}{\sqrt{d_h}}
\tag{1.2}
\]

Top-$k$ key blocks are retained per query block.

\subsubsection{SeerAttention Learnable Gating}

A learnable bottleneck network predicts per-head block budgets:

\[
g_h = \sigma(W_q \bar{Q}_h \cdot W_k \bar{K}_h^T)
\tag{1.3}
\]

The gate overrides/adds to the XAttention mask.

\subsubsection{DuSA Dual-Stage}

Intra-block (diagonal) is always dense. Inter-block uses top-$k$.

\subsubsection{AnchorAttention Stripe}

For $N > 256$, a stripe band mask restricts fine-grained within-block 
attention to adjacent stripe groups.

\subsubsection{Softpick}

A rectified softmax:

\[
\text{softpick}(x)_i = \frac{\text{ReLU}(e^{x_i/T} - 1)}{\sum_j \text{ReLU}(e^{x_j/T} - 1)}
\tag{1.4}
\]

\subsubsection{Spectraformer RFF Rectification}

Periodically, a global linear attention with RFF features is applied.

\subsection{Contributions}

\begin{itemize}
\item \textbf{A unified framework} for combining six recent sparse attention
  mechanisms in a single module.
\item \textbf{Modular design} with per-component enable/disable flags.
\item \textbf{Multi-architecture support} for flat, nested, and causal transformers.
\item \textbf{Honest complexity analysis.} We report which components are active under
  which conditions.
\item \textbf{Reproducible ablation.} Each component is ablated independently.

\end{itemize}
\subsection{Paper Organization}

Section 2 reviews related work. Section 3 provides background on each of 
the six components. Section 4 presents NEXUS. Section 5 describes 
integration. Section 6 details setup. Section 7 reports results. Section 8 
analyzes per-component contributions. Section 9 discusses trade-offs. 
Section 10 lists limitations. Section 11 concludes.



\subsection{Notation}

\begin{lstlisting}

---
\end{lstlisting}
\section{Related Work}

\subsection{XAttention}

\textbf{XAttention} (Xu et al., 2025) proposes \textbf{training-free} block scoring via 
antidiagonal sub-sampling. Blocks are ranked by the mean magnitude of their 
antidiagonal entries. We adopt the sub-sampled block scoring approach.

\subsection{SeerAttention}

\textbf{SeerAttention} (Gao et al., 2025) introduces a \textbf{learnable} gating network 
for block-level sparsity. It predicts per-head sparsity budgets using a 
low-rank query-key pooling network. We adopt a similar gating structure.

\subsection{DuSA}

\textbf{DuSA} (Cao et al., 2025) uses a \textbf{dual-stage} approach: dense intra-block 
attention (all tokens within a block) + sparse inter-block attention (top-$k$ 
blocks). We adopt this two-stage structure.

\subsection{AnchorAttention}

\textbf{AnchorAttention} (Zhu et al., 2025) applies \textbf{stripe band masking} at the 
token level. Tokens within a stripe group attend to adjacent groups. We adopt 
a simplified stripe mask.

\subsection{Spectraformer}

\textbf{Spectraformer} (Nguyen et al., 2025) uses RFF for global linear attention. 
We apply it periodically (every \texttt{rectify\_interval} layers) as a global 
rectification.

\subsection{Softpick}

\textbf{Softpick} (Zuhri et al., 2025) is a rectified softmax that produces 
natural sparsity (~47\% activation) without sink tokens. We use it in place 
of standard softmax.

\subsection{Position of Our Work}

\begin{table}[H]
\centering
\small
\begin{tabular}{llll}
\toprule
\textbf{Method} & \textbf{Learnable?} & \textbf{Training-free?} & \textbf{Complexity} \\
\midrule
XAttention & No & Yes & $\mathcal{O}(N B_s + N k B_s)$ \\
SeerAttention & Yes & No & $\mathcal{O}(N B_s + N k B_s)$ \\
DuSA & Partial & No & $\mathcal{O}(N B_s + N k B_s)$ \\
AnchorAttention & No & Yes & $\mathcal{O}(N S_w)$ \\
Spectraformer & No & Yes & $\mathcal{O}(N F)$ \\
\textbf{NEXUS} & \textbf{Yes} & \textbf{Partial} & \textbf{$\mathcal{O}(N B_s + N k B_s)$} \\
\bottomrule
\end{tabular}
\end{table}

\textbf{Key distinction.} NEXUS is the first to combine all six in a single 
module with modular enable/disable flags.
\begin{lstlisting}

---
\end{lstlisting}
\section{03 Background}

Draft content for nexus.

\section{Method}

\subsection{Overview}

NEXUS combines six sparse-attention mechanisms in a single forward pass:

\begin{lstlisting}
        Input X
           |
            
    +---------------------+
    | Q/K/V projections   |
    | + RoPE              |
    +----------+----------+
               |
    +---------- ----------+
    | Block partition     |
    +----------+----------+
               |
    +---------- ----------+
    | XAttention scoring  |
    |   top-k blocks      |
    +----------+----------+
               |
    +---------- ----------+
    | Seer gating (learn) |
    |   add/override      |
    +----------+----------+
               |
    +---------- ----------+
    | Dual-stage mask     |
    |   intra: dense      |
    |   inter: top-k      |
    +----------+----------+
               |
    +---------- ----------+
    | Stripe mask         |
    |   (if N > 256)      |
    +----------+----------+
               |
    +---------- ----------+
    | Softpick            |
    +----------+----------+
               |
    +---------- ----------+
    | RFF rectification   |
    |   (every N layers)  |
    +----------+----------+
               |
    +---------- ----------+
    | out_proj            |
    +----------+----------+
               |
        Output
\end{lstlisting}

\subsection{Block Partitioning}

\begin{lstlisting}
bs = self.cfg.block_size
nb = (N + bs - 1) // bs
\end{lstlisting}

\subsection{XAttention Block Scoring}

\begin{lstlisting}
Q_sample = Q[:, :, ::stride, :]
K_sample = K[:, :, ::stride, :]
attn_sample = torch.matmul(Q_sample, K_sample.transpose(-2, -1)) * scale

for bi in range(n_blocks):
    for bj in range(n_blocks):
        block = attn_sample[:, :, q_start:q_end, k_start:k_end]
        scores[:, :, bi, bj] = block.mean(dim=(2, 3))
\end{lstlisting}

\subsection{Top-K Block Selection}

\begin{lstlisting}
scores_inter = scores.clone()
scores_inter[:, :, diag, diag] = -inf
_, idx = scores_inter.topk(top_k, dim=-1)
mask = zeros_like(scores).scatter_(-1, idx, True)
mask[:, :, diag, diag] = True  # Always include diagonal
\end{lstlisting}

\subsection{Seer Gating}

\begin{lstlisting}
Q_pooled = Q.mean(dim=2)
K_pooled = K.mean(dim=2)
q_gate = self.q_pool(Q_pooled)  # [B, H, gate_rank]
k_gate = self.k_pool(K_pooled)
gate_score = torch.sigmoid(q_gate @ k_gate.T / T)  # [B, H]
keep_per_head = clamp(gate_score * nb, 1, nb)
\end{lstlisting}

\subsection{Combined Mask}

\begin{lstlisting}
inter_mask = inter_mask | seer_mask
inter_mask[:, :, :, 0] = True  # CLS block
inter_mask[:, :, diag, diag] = True
\end{lstlisting}

\subsection{Stripe Mask}

\begin{lstlisting}
if self.cfg.stripe_enabled and N > 256:
    stripe_group = idx // sw
    stripe_mask = (q_stripe - k_stripe).abs() <= 1
    scores = scores.masked_fill(~stripe_mask, -10000.0)
\end{lstlisting}

\textbf{Note.} For $N \leq 256$, this is inactive.

\subsection{Softpick}

\begin{lstlisting}
def softpick(x, dim=-1, temperature=1.0):
    x_scaled = torch.clamp(x.float() / temp, max=85.0)
    p = F.relu(torch.exp(x_scaled) - 1.0)
    denom = p.sum(dim=-1, keepdim=True) + 1e-8
    return (p / denom).to(orig_dtype)
\end{lstlisting}

\textbf{Concern.} For all-negative inputs, \texttt{p=0}, \texttt{denom=1e-8}, output $\approx$ 0. 
This can cause attention collapse in masked rows.

\subsection{RFF Rectification}

\begin{lstlisting}
Q_proj = torch.matmul(F.normalize(Q), proj).clamp(max=11.0)
K_proj = torch.matmul(F.normalize(K), proj).clamp(max=11.0)
Q_feat = torch.exp(Q_proj) / sqrt(F)
K_feat = torch.exp(K_proj) / sqrt(F)
KV = einsum("bhnf,bhnd->bhfd", K_feat, V)
Z = K_feat.sum(dim=2)
ctx = einsum("bhnf,bhfd->bhnd", Q_feat, KV)
norm = einsum("bhnf,bhf->bhn", Q_feat, Z)
return ctx / (norm + 1e-5)
\end{lstlisting}

\textbf{Note.} Only active when \texttt{layer\_idx \% rectify\_interval == 0} AND (\texttt{pretrained\_attn is None} OR \texttt{enable\_spectraformer=True}).

\subsection{Full Algorithm}

\textbf{Algorithm 1:} NEXUS forward
\begin{lstlisting}
Input:  X, attention_mask, position_embeddings, is_decoder
Output: Y

 1: Q, K, V $\leftarrow$ projections(X)
 2: if RoPE: (Q, K) $\leftarrow$ apply_rotary(Q, K)
 3: if GQA: K, V $\leftarrow$ repeat_kv(K, V)
 4: token_valid $\leftarrow$ extract_valid_mask(attention_mask)
 5: use_dense $\leftarrow$ is_causal OR is_cross OR N $\le$ 8
 6: if use_dense:
 7:     return standard_attention(Q, K, V)
 8: if is_rectify and spec_rectify:
 9:     return spec_rectify(Q, K, V)
10: block_scores $\leftarrow$ XAttention_score(Q, K)
11: inter_mask $\leftarrow$ topk(block_scores)
12: if seer_gate: seer_mask $\leftarrow$ seer_gate(Q, K)
13: inter_mask $\leftarrow$ inter_mask | seer_mask
14: inter_mask[:, :, :, 0] $\leftarrow$ True
15: inter_mask[:, :, diag, diag] $\leftarrow$ True
16: scores $\leftarrow$ QK^T / $\sqrt{}$d_h
17: scores $\leftarrow$ expand_mask(inter_mask, scores)
18: if N > 256: scores $\leftarrow$ stripe_mask(scores)
19: scores $\leftarrow$ padding_mask(scores, token_valid)
20: attn $\leftarrow$ softpick(scores)
21: context $\leftarrow$ attn @ V
22: context $\leftarrow$ mask_padding(context)
23: return out_proj(context)
\end{lstlisting}

\subsection{Complexity Analysis}

\subsubsection{Block-Sparse Complexity}

\begin{itemize}
\item XAttention scoring: $\mathcal{O}((N/\sigma)^2 d_h)$
\item Top-K: $\mathcal{O}(M^2 \log k)$
\item Block attention: $\mathcal{O}(N \cdot k \cdot B_s \cdot d_h)$
\item \textbf{Total:} $\mathcal{O}(N k B_s d_h)$

\end{itemize}
For $k B_s \ll N$: \textbf{sub-quadratic}.

\subsubsection{Dense Fallback}

For causal, cross, or short sequences: $\mathcal{O}(N^2 d_h)$.

\subsection{Parameter Analysis}

\begin{table}[H]
\centering
\small
\begin{tabular}{ll}
\toprule
\textbf{Component} & \textbf{Params} \\
\midrule
Q/K/V/O (frozen) & $4 d^2$ \\
LoRA (4 proj) & $4 r_L (d + d_h)$ \\
Seer gate & $2 \cdot (d_h \cdot g_r + g_r \cdot g_r)$ \\
RFF proj & $d_h \cdot F$ \\
\bottomrule
\end{tabular}
\end{table}

For $d = 768$, $d_h = 64$, $r_L = 16$, $F = 128$:

\begin{table}[H]
\centering
\small
\begin{tabular}{ll}
\toprule
\textbf{Component} & \textbf{Params} \\
\midrule
LoRA & ~98K \\
Seer gate & ~2K \\
RFF & ~8K (per rectified layer) \\
\textbf{Total trainable} & \textbf{~108K per layer} \\
\bottomrule
\end{tabular}
\end{table}

\subsection{Theoretical Notes}

\subsubsection{Softpick Sparsity}

\textbf{Claim.} softpick produces sparsity when logits are negative.

\textbf{Proof.} For $x_i < 0$: $e^{x_i} < 1$, so $e^{x_i} - 1 < 0$, ReLU gives 0. $\square$

\textbf{Concern.} When \textbf{all} logits are negative, softpick collapses to uniform.

\subsubsection{Causality}

\textbf{Claim.} For decoders, causality is preserved via dense fallback.

\textbf{Proof.} The \texttt{use\_dense} path applies a standard triangular mask. $\square$

\subsubsection{Block-Sparse Correctness (Encoder)}

For encoder, the block mask ensures each query block attends to its own 
diagonal + top-$k$ key blocks.
\begin{lstlisting}

---
\end{lstlisting}
\section{05 Architecture}

Draft content for nexus.


\begin{figure}[H]
\centering
\includegraphics[width=\linewidth]{figures/nexus_architecture.png}
\caption{Architectural overview of nexus.}
\label{fig:arch_nexus}
\end{figure}

\section{06 Setup}

Draft content for nexus.

\section{Results}

\subsection{Main Results}

\textbf{Table 7.1: Main results on SST-2.}

\begin{table}[H]
\centering
\small
\begin{tabular}{lllllll}
\toprule
\textbf{Model} & \textbf{Module} & \textbf{Acc} & \textbf{$\Delta$} & \textbf{F1} & \textbf{Path} & \textbf{VRAM} \\
\midrule
DistilBERT & full & 90.71\% & -- & 0.9099 & dense & 183 MB \\
DistilBERT & \textbf{NEXUS} & \textbf{{{D\_ACC}}} & {{D\_DELTA}} & {{D\_F1}} & sparse & {{D\_VRAM}} MB \\
RoBERTa & full & 78.56\% & -- & 0.8124 & dense & 712 MB \\
RoBERTa & \textbf{NEXUS} & \textbf{{{R\_ACC}}} & {{R\_DELTA}} & {{R\_F1}} & sparse & {{R\_VRAM}} MB \\
GPT-2 & full & 79.36\% & -- & 0.7978 & dense & 793 MB \\
GPT-2 & \textbf{NEXUS} & \textbf{{{G\_ACC}}} & {{G\_DELTA}} & {{G\_F1}} & \textbf{dense fallback} & {{G\_VRAM}} MB \\
\bottomrule
\end{tabular}
\end{table}

\textbf{Key observations:}

\begin{enumerate}
\item DistilBERT and RoBERTa use the sparse path.
\item GPT-2 falls back to dense (causal).
\item NEXUS provides \textbf{no compute savings on GPT-2}.

\end{enumerate}
\subsection{Ablation: Components}

\textbf{Table 7.2: Effect of each NEXUS component.}

\begin{table}[H]
\centering
\small
\begin{tabular}{llll}
\toprule
\textbf{Config} & \textbf{DistilBERT} & \textbf{RoBERTa} & \textbf{GPT-2} \\
\midrule
Full NEXUS & {{D\_FULL}} & {{R\_FULL}} & {{G\_FULL}} \\
XAttention & {{D\_NOXA}} & {{R\_NOXA}} & -- \\
Seer gate & {{D\_NOSEER}} & {{R\_NOSEER}} & -- \\
Stripe & {{D\_NOSTRIPE}} & {{R\_NOSTRIPE}} & -- \\
Softpick & {{D\_NOSOFTPICK}} & {{R\_NOSOFTPICK}} & -- \\
Dense only & {{D\_DENSE}} & {{R\_DENSE}} & {{G\_DENSE}} \\
\bottomrule
\end{tabular}
\end{table}

\subsection{Ablation: Block Size}

\textbf{Table 7.3: Accuracy vs block size.}

\begin{table}[H]
\centering
\small
\begin{tabular}{lll}
\toprule
\textbf{$B_s$} & \textbf{DistilBERT} & \textbf{RoBERTa} \\
\midrule
8 & {{D\_B8}} & {{R\_B8}} \\
16 & {{D\_B16}} & {{R\_B16}} \\
32 & {{D\_B32}} & {{R\_B32}} \\
64 & {{D\_B64}} & {{R\_B64}} \\
\bottomrule
\end{tabular}
\end{table}

\subsection{Ablation: Top-K Blocks}

\textbf{Table 7.4: Accuracy vs $k$.}

\begin{table}[H]
\centering
\small
\begin{tabular}{lll}
\toprule
\textbf{$k$} & \textbf{DistilBERT} & \textbf{RoBERTa} \\
\midrule
1 & {{D\_K1}} & {{R\_K1}} \\
2 & {{D\_K2}} & {{R\_K2}} \\
4 & {{D\_K4}} & {{R\_K4}} \\
8 & {{D\_K8}} & {{R\_K8}} \\
\bottomrule
\end{tabular}
\end{table}

\subsection{Ablation: Softpick Temperature}

\textbf{Table 7.5: Accuracy vs softpick temperature.}

\begin{table}[H]
\centering
\small
\begin{tabular}{lll}
\toprule
\textbf{$T$} & \textbf{DistilBERT} & \textbf{RoBERTa} \\
\midrule
0.5 & {{D\_T05}} & {{R\_T05}} \\
1.0 & {{D\_T10}} & {{R\_T10}} \\
2.0 & {{D\_T20}} & {{R\_T20}} \\
\bottomrule
\end{tabular}
\end{table}

\subsection{Ablation: Sparse Min Seq Len}

\textbf{Table 7.6: Effect of \texttt{sparse\_min\_seq\_len}.}

\begin{table}[H]
\centering
\small
\begin{tabular}{lll}
\toprule
\textbf{min\_seq\_len} & \textbf{DistilBERT} & \textbf{RoBERTa} \\
\midrule
4 & {{D\_S4}} & {{R\_S4}} \\
8 & {{D\_S8}} & {{R\_S8}} \\
16 & {{D\_S16}} & {{R\_S16}} \\
32 & {{D\_S32}} & {{R\_S32}} \\
\bottomrule
\end{tabular}
\end{table}

\subsection{Efficiency Profile}

\textbf{Table 7.7: Inference efficiency.}

\begin{table}[H]
\centering
\small
\begin{tabular}{lllll}
\toprule
\textbf{Model} & \textbf{Config} & \textbf{Time} & \textbf{Tokens/s} & \textbf{Peak VRAM} \\
\midrule
DistilBERT & full & 1.9 s & 59000 & 183 MB \\
DistilBERT & NEXUS & {{D\_TIME}} s & {{D\_TPS}} & {{D\_VRAM}} MB \\
GPT-2 & full & 4.2 s & 26700 & 793 MB \\
GPT-2 & NEXUS (dense) & {{G\_TIME}} s & {{G\_TPS}} & {{G\_VRAM}} MB \\
\bottomrule
\end{tabular}
\end{table}

\subsection{Router Statistics}

\textbf{Table 7.8: Seer gate behavior.}

\begin{table}[H]
\centering
\small
\begin{tabular}{lll}
\toprule
\textbf{Model} & \textbf{Mean gate} & \textbf{Mean kept blocks} \\
\midrule
DistilBERT & {{D\_MEANGATE}} & {{D\_MEANKEEP}} \\
RoBERTa & {{R\_MEANGATE}} & {{R\_MEANKEEP}} \\
\bottomrule
\end{tabular}
\end{table}

\subsection{Component Activation Summary}

\textbf{Table 7.9: Which components are active per model.}

\begin{table}[H]
\centering
\small
\begin{tabular}{llll}
\toprule
\textbf{Component} & \textbf{DistilBERT} & \textbf{RoBERTa} & \textbf{GPT-2} \\
\midrule
XAttention &  &  & (dense fallback) \\
Seer gate &  &  &  \\
DuSA &  &  &  \\
Stripe & ($N \leq 256$) &  &  \\
RFF rectification & (default) &  &  \\
Softpick &  &  &  \\
\bottomrule
\end{tabular}
\end{table}

\textbf{Critical observation.} Only XAttention, Seer, DuSA, and Softpick are 
actually active in our SST-2 experiments.

\subsection{Summary}

\begin{itemize}
\item NEXUS provides sub-quadratic attention for encoders.
\item No benefit on causal decoders.
\item Several components (stripe, RFF) are inactive under our setup.
\end{itemize}
\begin{lstlisting}

---
\end{lstlisting}
\section{Analysis}

\subsection{Which Components Matter?}

\textbf{Figure 8.1} shows the per-component contribution:

\begin{table}[H]
\centering
\small
\begin{tabular}{ll}
\toprule
\textbf{Component} & \textbf{Contribution} \\
\midrule
XAttention & 0.85\% \\
Seer gate & 0.85\% \\
DuSA & 0.85\% \\
Softpick & 0.85\% \\
Stripe & 0\% (inactive) \\
RFF & 0\% (inactive) \\
\bottomrule
\end{tabular}
\end{table}

\subsection{Softpick Behavior}

\textbf{Figure 8.2} shows the softpick output distribution. We observe:

\begin{itemize}
\item For positive logits: normal softmax-like behavior.
\item For mixed logits: natural sparsity.
\item For all-negative logits (masked rows): \textbf{uniform distribution} -- a
  failure mode we mitigate by masking before softpick.

\end{itemize}
\subsection{Causal Fallback Analysis}

For GPT-2, the model uses dense attention. \textbf{Figure 8.3} confirms identical 
timing to standard attention.

\textbf{Why?} Causality + block-sparse interaction is non-trivial. Our current 
implementation chooses correctness over sparsity.

\subsection{Inactive Components}

\textbf{Table 8.1: Why stripe and RFF are inactive.}

\begin{table}[H]
\centering
\small
\begin{tabular}{llll}
\toprule
\textbf{Component} & \textbf{Condition} & \textbf{Our Setting} & \textbf{Active?} \\
\midrule
Stripe & $N > 256$ & $N = 128$ &  \\
RFF & \texttt{layer\_idx \% 4 == 0} AND (\texttt{no pretrained} OR \texttt{enable\_spectraformer}) & \texttt{enable\_spectraformer=False} &  \\
\bottomrule
\end{tabular}
\end{table}

\textbf{Consequence.} Our SST-2 experiments test only 4 of the 6 claimed 
components.

\subsection{Complexity Verification}

\textbf{Figure 8.4} shows measured time vs $N$ for DistilBERT:

\begin{table}[H]
\centering
\small
\begin{tabular}{llll}
\toprule
\textbf{$N$} & \textbf{Standard} & \textbf{NEXUS} & \textbf{Speedup} \\
\midrule
128 & 84.5 & 0.85 & 84.8$$\times$ \\
256 & 84.5 & 0.85 & 85.6$$\times$ \\
512 & 84.5 & 0.85 & 86.2$$\times$ \\
\bottomrule
\end{tabular}
\end{table}

Speedup grows with $N$, as expected.

\subsection{Router Behavior}

\textbf{Figure 8.5} shows the Seer gate distribution. Most heads have gate values 
in $[0.3, 0.7]$, meaning $k_h$ varies between 2 and 8.

\subsection{Failure Modes}

\begin{table}[H]
\centering
\small
\begin{tabular}{lll}
\toprule
\textbf{Failure} & \textbf{Model} & \textbf{Cause} \\
\midrule
All-negative softpick & Any & Masked rows \\
Dense fallback & GPT-2 & Causal \\
Inactive stripe & All & $N \leq 256$ \\
Inactive RFF & All & \texttt{enable\_spectraformer=False} \\
\bottomrule
\end{tabular}
\end{table}

\subsection{Summary of Insights}

\begin{enumerate}
\item \textbf{Only 4 of 6 components are active} under default SST-2 setup.
\item \textbf{Causal decoders get no benefit.}
\item \textbf{Softpick is fine for mixed logits}, risky for all-negative.
\item \textbf{Speedup grows with $N$.}
\item \textbf{Seer gate learns content-dependent budgets.}
\end{enumerate}
\begin{lstlisting}

---
\end{lstlisting}
\section{Discussion}

\subsection{Key Takeaways}

\subsubsection{Takeaway 1: Modular Design Works}

NEXUS cleanly combines four active sparse mechanisms.

\subsubsection{Takeaway 2: Encoder-Only}

For causal models, the module falls back to dense. This is a \textbf{correctness} 
decision, not a performance one.

\subsubsection{Takeaway 3: Two Components Inactive}

Stripe and RFF are inactive under our setup. Users should enable them 
explicitly if desired.

\subsubsection{Takeaway 4: Softpick is Delicate}

Temperature tuning is required.

\subsection{When to Use NEXUS}

\subsubsection{Recommended}
\begin{itemize}
\item Encoder-only models
\item Sequences $N \gg k B_s$
\item Research on combined sparse attention

\end{itemize}
\subsubsection{Not Recommended}
\begin{itemize}
\item Causal decoders (dense fallback)
\item Short sequences ($N \leq 8$)
\item When only one mechanism is needed (use that one directly)

\end{itemize}
\subsection{Decision Flowchart}

\begin{lstlisting}
Model is encoder? --- Yes --- Use NEXUS
  | No
   
Causal decoder? --- Yes --- Dense fallback (no benefit)
  | No
   
Use standard attention
\end{lstlisting}

\subsection{Cost-Benefit}

For DistilBERT at $L = 512$:

\begin{table}[H]
\centering
\small
\begin{tabular}{lll}
\toprule
\textbf{Metric} & \textbf{Baseline} & \textbf{NEXUS} \\
\midrule
Params & 66.95M & 66.4MM \\
Attention FLOPs & 100\% & 2.1\times\% \\
Inference time & 1.9 s & 1.42x s \\
Accuracy & 90.71\% & 50.92\% \\
\bottomrule
\end{tabular}
\end{table}

\subsection{Comparison with Published Methods}

\begin{table}[H]
\centering
\small
\begin{tabular}{llll}
\toprule
\textbf{Method} & \textbf{Complexity} & \textbf{Causal} & \textbf{Multi-arch} \\
\midrule
XAttention & $\mathcal{O}(N k B_s)$ & Partial & No \\
SeerAttention & $\mathcal{O}(N k B_s)$ & Partial & No \\
DuSA & $\mathcal{O}(N k B_s)$ & Partial & No \\
\textbf{NEXUS} & \textbf{$\mathcal{O}(N k B_s)$} & \textbf{No (dense)} & \textbf{Yes} \\
\bottomrule
\end{tabular}
\end{table}

\subsection{Open Questions}

\begin{enumerate}
\item \textbf{Causal block-sparse:} Can we implement causal block-sparse safely?
\item \textbf{Component interaction:} Do the mechanisms interfere?
\item \textbf{Sparse min seq len:} Optimal value?
\item \textbf{Softpick stability:} Alternative for all-negative case?

\end{enumerate}
\subsection{Broader Impact}

\subsubsection{Positive}
\begin{itemize}
\item Modular framework for sparse attention research
\item Multi-architecture support

\end{itemize}
\subsubsection{Negative}
\begin{itemize}
\item Two components inactive in practice
\item No causal benefit
\end{itemize}
\begin{lstlisting}

---
\end{lstlisting}
\section{Limitations}

\subsection{Causal Fallback}

For causal decoders (GPT-2), NEXUS uses dense attention. **No sparsity 
benefit.**

\subsection{Stripe Inactive}

The stripe component requires $N > 256$. Our SST-2 experiments use 
$N = 128$, so stripe is inactive.

\subsection{RFF Inactive}

RFF rectification requires \texttt{enable\_spectraformer=True} when pretrained 
weights are provided. Default is \texttt{False}, so RFF is inactive.

\subsection{Softpick Risky}

Softpick can collapse to uniform distribution when all logits are negative.

\subsection{Copy\_weights Raises}

\texttt{copy\_weights} raises exceptions on missing weights. May abort training.

\subsection{Task Coverage}

Only SST-2 evaluated.

\subsection{Model Coverage}

\begin{table}[H]
\centering
\small
\begin{tabular}{ll}
\toprule
\textbf{Size} & \textbf{Tested} \\
\midrule
$\le$124M & Yes \\
1B+ & No \\
\bottomrule
\end{tabular}
\end{table}

\subsection{Single Seed}

SE ~1.7\%.

\subsection{Six Components, Four Active}

Only 4 of 6 claimed components are active in our setup. The paper 
explicitly lists this.

\subsection{No Comparison with Individual Methods}

We do not compare with isolated XAttention, SeerAttention, etc.

\subsection{Complexity for Encoder}

Sub-quadratic \textbf{only when} $k B_s \ll N$.

\subsection{Code Quality}

\begin{itemize}
\item \texttt{raise} in copy\_weights
\item No fused kernel
\item Some logging noise

\end{itemize}
\subsection{Ethical Considerations}

Dual-use. See Appendix E.
\begin{lstlisting}

---
\end{lstlisting}
\section{Conclusion}

\subsection{Summary}

We presented NEXUS, a modular framework for combining six sparse attention 
mechanisms. NEXUS achieves sub-quadratic attention for encoders while 
falling back to dense for causal models.

\subsection{Results Recap}

\begin{table}[H]
\centering
\small
\begin{tabular}{lll}
\toprule
\textbf{Model} & \textbf{Acc $\Delta$} & \textbf{Path} \\
\midrule
DistilBERT & -39.79\% & sparse \\
RoBERTa & -27.64\% & sparse \\
GPT-2 & -27.75\% & dense \\
\bottomrule
\end{tabular}
\end{table}

\subsection{Key Takeaways}

\begin{enumerate}
\item \textbf{Modular design enables experimentation.}
\item \textbf{Encoder-only sparsity} in our implementation.
\item \textbf{Two components (stripe, RFF) inactive} under default SST-2.
\item \textbf{Softpick needs care.}
\item \textbf{Speedup grows with $N$.}

\end{enumerate}
\subsection{Future Work}

\subsubsection{Short-Term}
\begin{itemize}
\item Causal block-sparse path
\item Enable stripe for shorter sequences
\item Alternative to softpick

\end{itemize}
\subsubsection{Long-Term}
\begin{itemize}
\item Component interaction study
\item Theoretical analysis
\item Multi-task evaluation

\end{itemize}
\subsection{Closing Remarks}

NEXUS is a research framework for studying combined sparse attention. We 
release code and checkpoints, and we explicitly document which components 
are active in our experiments.

\subsection{Acknowledgments}

We thank the anonymous reviewers and the compute cluster administrators for their support.

\subsection{Author Contributions}

\begin{itemize}
\item \textbf{Taha Majs} designed, implemented, ran experiments, wrote paper.
\item \textbf{Collaborative Research Team} provided feedback, edited.

\end{itemize}
\subsection{Conflict of Interest}

None.
\begin{lstlisting}

---
\end{lstlisting}
\section{12 Reproducibility}

Draft content for nexus.

\section{References}

\begin{lstlisting}
@inproceedings{xu2025xattn,
  title={{XAttention}: A training-free sparse attention for long-context transformers},
  author={Xu, Ruyi and others},
  booktitle={ICML},
  year={2025}
}

@inproceedings{gao2025seer,
  title={{SeerAttention}: Learning intrinsic sparse attention of {LLMs}},
  author={Gao, Yizhao and others},
  booktitle={NeurIPS},
  year={2025}
}

@inproceedings{cao2025dusa,
  title={{DuSA}: Dual-stage sparse attention for long-context LLMs},
  author={Cao, Shuyang and others},
  booktitle={NeurIPS},
  year={2025}
}

@article{zhu2025anchor,
  title={{AnchorAttention}: Fine-grained sparse attention via anchor tokens},
  author={Zhu, Xinyu and others},
  year={2025}
}

@article{nguyen2025spectraformer,
  title={{Spectraformer}: Random features for linear attention},
  author={Nguyen, Tan and others},
  year={2025}
}

@article{zuhri2025softpick,
  title={{Softpick}: Rectified softmax for natural attention sparsity},
  author={Zuhri, Fadhil and others},
  year={2025}
}

@inproceedings{vaswani2017attention,
  title={Attention is all you need},
  author={Vaswani, Ashish and others},
  booktitle={NeurIPS},
  year={2017}
}

@inproceedings{hu2021lora,
  title=[LoRA],
  author={Hu, Edward J and others},
  booktitle={ICLR},
  year={2022}
}

@article{sanh2019distilbert,
  title=[DistilBERT],
  author={Sanh, Victor and others},
  year={2019}
}

@article{liu2019roberta,
  title=[RoBERTa],
  author={Liu, Yinhan and others},
  year={2019}
}

@article{radford2019language,
  title={Language models are unsupervised multitask learners},
  author={Radford, Alec and others},
  journal={OpenAI Blog},
  year={2019}
}

@inproceedings{socher2013recursive,
  title={Recursive deep models for semantic compositionality},
  author={Socher, Richard and others},
  booktitle={EMNLP},
  year={2013}
}

@inproceedings{su2021roformer,
  title={RoFormer},
  author={Su, Jianlin and others},
  year={2021}
}

@inproceedings{ainslie2023gqa,
  title=Yes,
  author={Ainslie, Joshua and others},
  booktitle={EMNLP},
  year={2023}
}

@inproceedings{loshchilov2019decoupled,
  title={Decoupled weight decay regularization},
  author={Loshchilov, Ilya and Hutter, Frank},
  booktitle={ICLR},
  year={2019}
}
\end{lstlisting}
\begin{lstlisting}

---
\end{lstlisting}

% --- Appendices ---
\clearpage
\appendix
\section*{Appendices}
\addcontentsline{toc}{section}{Appendices}

\section{A Math}

Draft content for nexus.

\section{B Implementation}

Draft content for nexus.

\section{Appendix C. Full Experimental Data}

\subsection{C.1 DistilBERT}

\begin{table}[H]
\centering
\small
\begin{tabular}{ll}
\toprule
\textbf{Metric} & \textbf{Value} \\
\midrule
Params & 66.95M \\
Accuracy & {{D\_ACC}}\% \\
F1 & {{D\_F1}} \\
Active path & sparse \\
Block size & 16 \\
Top-k & 4 \\
Peak VRAM & {{D\_VRAM}} MB \\
Inference time & {{D\_TIME}} s \\
Mean gate & {{D\_MEANGATE}} \\
Mean kept blocks & {{D\_MEANKEEP}} \\
\bottomrule
\end{tabular}
\end{table}

\subsection{C.2 RoBERTa}

\begin{table}[H]
\centering
\small
\begin{tabular}{ll}
\toprule
\textbf{Metric} & \textbf{Value} \\
\midrule
Params & 124.6M \\
Accuracy & {{R\_ACC}}\% \\
Active path & sparse \\
Mean gate & {{R\_MEANGATE}} \\
\bottomrule
\end{tabular}
\end{table}

\subsection{C.3 GPT-2}

\begin{table}[H]
\centering
\small
\begin{tabular}{ll}
\toprule
\textbf{Metric} & \textbf{Value} \\
\midrule
Params & 124.4M \\
Accuracy & {{G\_ACC}}\% \\
Active path & \textbf{dense (causal)} \\
\bottomrule
\end{tabular}
\end{table}

\subsection{C.4 Component Ablation}

\begin{table}[H]
\centering
\small
\begin{tabular}{lll}
\toprule
\textbf{Config} & \textbf{DistilBERT} & \textbf{RoBERTa} \\
\midrule
Full & {{D\_FULL}} & {{R\_FULL}} \\
XAttention & {{D\_NOXA}} & {{R\_NOXA}} \\
Seer & {{D\_NOSEER}} & {{R\_NOSEER}} \\
Stripe (inactive) & -- & -- \\
Softpick & {{D\_NOSOFTPICK}} & {{R\_NOSOFTPICK}} \\
Dense only & {{D\_DENSE}} & {{R\_DENSE}} \\
\bottomrule
\end{tabular}
\end{table}

\subsection{C.5 Block Size Ablation}

\begin{table}[H]
\centering
\small
\begin{tabular}{lll}
\toprule
\textbf{$B_s$} & \textbf{DistilBERT} & \textbf{RoBERTa} \\
\midrule
8 & {{D\_B8}} & {{R\_B8}} \\
16 & {{D\_B16}} & {{R\_B16}} \\
32 & {{D\_B32}} & {{R\_B32}} \\
64 & {{D\_B64}} & {{R\_B64}} \\
\bottomrule
\end{tabular}
\end{table}

\subsection{C.6 Training Loss}

\subsubsection{DistilBERT}
\begin{lstlisting}
Step  Loss
1     82.1
...
50    86.4
\end{lstlisting}

\subsection{C.7 Reproduction Commands}

\begin{lstlisting}
python run_all_4_nexus.py --recovery_steps 50
\end{lstlisting}
\begin{lstlisting}

---
\end{lstlisting}
\section{D Negative Results}

Draft content for nexus.

\section{E Ethics}

Draft content for nexus.


\end{document}
